%!TEX root = ../../note.tex
\subsection{The Limit of a Sequence}

\subsubsection{Definition and first examples}
A \term{sequence} of real numbers is a function $X:\N\to\R$. We write
$x_n=X(n)$ for its $n$th term and denote the sequence by
$\set{x_n}_{n=1}^{\infty}$ or simply $\set{x_n}$.

\begin{definition}\label{def:limit}
    A sequence $\set{x_n}$ \term{converges} to $a\in\R$, written
    $\lim_{n\to\infty}x_n=a$ or $x_n\to a$, if
    \begin{equation*}
        (\forall \epsilon>0)(\exists n_\epsilon\in\N)(\forall n\geq n_\epsilon)
        \,\bigl[\abs{x_n-a}<\epsilon\bigr].
    \end{equation*}
\end{definition}

In words: for every tolerance $\epsilon>0$, all sufficiently late terms lie
within $\epsilon$ of $a$. The index $n_\epsilon$ is allowed to depend on
$\epsilon$ (smaller tolerances usually require later indices), and finitely
many early terms may lie anywhere. In the figure, $x_n=1/n$, $a=0$,
$\epsilon=0.3$: every term from $n_\epsilon=4$ on lies in the shaded band.

\begin{center}
\begin{tikzpicture}[x=0.55cm,y=2.4cm,>=Stealth]
    \fill[red!8] (0,0) rectangle (11,0.3);
    \draw[densely dashed,red!75!black] (0,0.3) -- (11,0.3) node[right] {$a+\epsilon$};
    \draw[->] (0,0) -- (11.5,0) node[right] {$n$};
    \draw[->] (0,0) -- (0,1.15) node[above] {$x_n$};
    \node[left] at (0,0) {$a=0$};
    \draw[densely dashed,blue!70!black] (4,0) -- (4,0.95);
    \node[below] at (4,0) {$n_\epsilon=4$};
    \foreach \n in {1,...,10} {
        \fill[teal!70!black] (\n,{1/\n}) circle (1.7pt);
    }
    \node[teal!70!black,anchor=west] at (1.3,0.85) {$x_n=1/n$};
\end{tikzpicture}
\end{center}

The interval
\[
    V_\epsilon(a)=\set{x\in\R:\abs{x-a}<\epsilon}=(a-\epsilon,a+\epsilon)
\]
is called the \term{$\epsilon$-neighbourhood} of $a$. Thus $x_n\to a$ if and
only if every $\epsilon$-neighbourhood of $a$ contains all terms of the
sequence from some index on.

To prove $x_n\to a$ from the definition, one fixes $\epsilon>0$, estimates
$\abs{x_n-a}$ from above by something simple, and chooses $n_\epsilon$ ---
usually by the Archimedean property --- to make that bound less than
$\epsilon$.

\begin{example}\label{ex:1/n}
    $\lim_{n\to\infty}\frac1n=0$.
\end{example}

\begin{proof}
    Given $\epsilon>0$, choose $n_\epsilon$ with $1/n_\epsilon<\epsilon$
    (Corollary~\ref{cor:archimedes}(b)). For $n\geq n_\epsilon$,
    $\abs{\frac1n-0}=\frac1n\leq\frac1{n_\epsilon}<\epsilon$.
\end{proof}

\begin{example}
    $\lim_{n\to\infty}\frac{4n+3}{n+1}=4$.
\end{example}

\begin{proof}
    Given $\epsilon>0$, choose $n_\epsilon$ with $1/n_\epsilon<\epsilon$. For
    $n\geq n_\epsilon$,
    \begin{equation*}
        \abs{\frac{4n+3}{n+1}-4}
        =\frac1{n+1}<\frac1n\leq\frac1{n_\epsilon}<\epsilon.
    \end{equation*}
\end{proof}

\subsubsection{Uniqueness of the limit}
\begin{proposition}\label{prop:eps-equal}
    If $x,y\in\R$ and $\abs{x-y}<\epsilon$ for every $\epsilon>0$, then $x=y$.
\end{proposition}

\begin{proof}
    If $x\ne y$, take $\epsilon=\abs{x-y}/2>0$; then $\abs{x-y}>\epsilon$.
\end{proof}

\begin{theorem}\label{thm:limit-unique}
    A sequence has at most one limit.
\end{theorem}

If $x_n\to a$ and $x_n\to b$, a late enough term is close to both $a$ and
$b$, so by the triangle inequality $a$ and $b$ are close to each other.

\begin{proof}
    Let $\epsilon>0$. Choose $n_1$ with $\abs{x_n-a}<\epsilon/2$ for
    $n\geq n_1$ and $n_2$ with $\abs{x_n-b}<\epsilon/2$ for $n\geq n_2$. For
    $n\geq\max\set{n_1,n_2}$,
    \begin{equation*}
        \abs{a-b}\leq\abs{a-x_n}+\abs{x_n-b}<\frac\epsilon2+\frac\epsilon2=\epsilon.
    \end{equation*}
    Since $\epsilon>0$ was arbitrary, $a=b$ by
    Proposition~\ref{prop:eps-equal}.
\end{proof}

\subsubsection{Divergence and boundedness}
A sequence that does not converge to any real number is said to
\term{diverge}.

\begin{example}\label{ex:0101}
    The sequence $0,1,0,1,\ldots$ ($x_n=0$ for odd $n$, $x_n=1$ for even $n$)
    diverges: both values recur arbitrarily late, so the terms cannot all
    stay near one number.
\end{example}

\begin{proof}
    Suppose $x_n\to a$ and take $\epsilon=1/2$. Choosing an odd and an even
    index $n\geq n_\epsilon$ gives $\abs{0-a}<\frac12$ and $\abs{1-a}<\frac12$,
    i.e.\ $-\frac12<a<\frac12$ and $\frac12<a<\frac32$, which is impossible.
\end{proof}

\begin{definition}
    A sequence $\set{x_n}$ is \term{bounded} if there is $M\in\R$ with
    $\abs{x_n}\leq M$ for all $n\in\N$.
\end{definition}

\begin{theorem}\label{thm:conv-bounded}
    Every convergent sequence is bounded.
\end{theorem}

Convergence controls all terms from some index on; only finitely many terms
remain, and a finite set is bounded.

\begin{proof}
    Let $x_n\to a$. With $\epsilon=1$, choose $n_1$ such that
    $\abs{x_n-a}<1$ for $n\geq n_1$; then
    $\abs{x_n}\leq\abs{x_n-a}+\abs a<1+\abs a$ for these $n$. Hence
    \begin{equation*}
        M=\max\set{\abs{x_1},\ldots,\abs{x_{n_1-1}},\,1+\abs{a}}
    \end{equation*}
    bounds every term (if $n_1=1$, take $M=1+\abs a$).
\end{proof}

The converse is false: the sequence of Example~\ref{ex:0101} is bounded but
divergent. On the other hand, by Theorem~\ref{thm:conv-bounded} an
\emph{unbounded} sequence must diverge.

\begin{example}
    The following sequences are unbounded, hence divergent.

    (a) $a_n=\frac{n^2+1}{n}$. Here $\abs{a_n}=n+\frac1n$. Given $M>0$,
    choose $n_M\in\N$ with $n_M>M$; then $\abs{a_{n_M}}>M$.

    (b) $a_n=(-1)^n n$. Here $\abs{a_n}=n$, and again $\abs{a_{n_M}}>M$
    whenever $n_M>M$.
\end{example}

\subsubsection{Infinite limits}
\begin{definition}
    A sequence $\set{x_n}$ \term{diverges to $+\infty$}, written
    $\lim_{n\to\infty}x_n=+\infty$, if for every $M>0$ there is
    $n_M\in\N$ such that $x_n>M$ for all $n\geq n_M$. It diverges to
    $-\infty$ if for every $M>0$ there is $n_M$ such that $x_n<-M$ for all
    $n\geq n_M$.
\end{definition}

Such a sequence is unbounded and therefore \emph{divergent}; the notation
$\lim x_n=\pm\infty$ only records the manner in which it diverges.

\begin{example}
    $\lim_{n\to\infty}\frac{n^2}{1-n}=-\infty$.
\end{example}

\begin{proof}
    For $n\geq2$ we have $0<n-1<n$, so
    $\frac{n^2}{n-1}>\frac{n^2}{n}=n$, i.e.\ $\frac{n^2}{1-n}<-n$. Given
    $M>0$, choose $n_M\in\N$ with $n_M\geq2$ and $n_M>M$. For $n\geq n_M$,
    \begin{equation*}
        \frac{n^2}{1-n}<-n\leq-n_M<-M.
    \end{equation*}
\end{proof}

\begin{example}
    $\lim_{n\to\infty}\paren{\sqrt{n}+1}=+\infty$.
\end{example}

\begin{proof}
    Given $M>0$, we must make $\sqrt n>M-1$ for all late $n$. Choose
    $n_M\in\N$ with $n_M>M^2$. For $n\geq n_M$,
    $\sqrt n\geq\sqrt{n_M}>M>M-1$, so $\sqrt n+1>M$.
\end{proof}

\subsubsection{Limits and arithmetic operations}
% The lecture stops after the first two parts; the theorem is continued
% (sums, products, quotients of limits) in the next lecture.
\begin{theorem}\label{thm:limit-laws}
    Let $\lim_{n\to\infty} x_n = a$ and $c\in\R$. Then
    \begin{enumerate}[label=(\arabic*)]
        \item $\lim_{n \to \infty} c = c$ (the constant sequence);
        \item $\lim_{n \to \infty} cx_n = ca$.
    \end{enumerate}
\end{theorem}

\begin{proof}
    (1) $\abs{c-c}=0<\epsilon$ for every $n$.

    (2) If $c=0$, this is (1). If $c\neq0$, given $\epsilon>0$ choose
    $n_\epsilon$ with $\abs{x_n-a}<\epsilon/\abs c$ for $n\geq n_\epsilon$;
    then $\abs{cx_n-ca}=\abs c\abs{x_n-a}<\epsilon$.
\end{proof}
